\section{Equations, waves and what is counted}\label{sec:setting}

This section fixes the normalizations and records the form of the waves
that are counted.

\subsection{From the evolution equation to the profile equation}

For a travelling wave \(u=w(\xi)\), \(\xi=x-ct\), one integration turns
\eqref{eq:pde} into \eqref{eq:integrated}:
\begin{align}
 &u_t+uu_x+\sum_{j=1}^{p}b_j\partial_x^{j+1}u=0,\qquad b_p\ne0,
 \label{eq:pde}\\
 &\sum_{j=1}^{p}b_jw^{(j)}-cw+\frac12w^2+K=0,\label{eq:integrated}
\end{align}
with a constant \(K\). Let \(\kappa\ne0\), and let \(w_\ast\) be a
constant solution of \eqref{eq:integrated}, a root of \(w^2/2-cw+K\).
The substitution \eqref{eq:scaling} turns \eqref{eq:integrated} into the
profile equation \eqref{eq:main} with coefficients
\eqref{eq:coefficients}:
\begin{align}
 w(\xi)&=w_\ast+b_p\kappa^p\,v(\kappa\xi),\label{eq:scaling}\\
 a_j&=\frac{b_j}{b_p}\kappa^{j-p}\quad(1\le j\le p-1),\qquad
 a_0=\frac{w_\ast-c}{b_p\kappa^p}.\label{eq:coefficients}
\end{align}
Thus \(a_1,\ldots,a_{p-1}\) record the equation, up to the scale
\(\kappa\), and \(a_0\) records the speed of the wave relative to the
background \(w_\ast\). The scale is free at this stage. The period of
the wave will fix it: if \(v\) is rational in \(e^{z}\), then \(w\) is
rational in \(e^{\kappa\xi}\), so that \(\kappa\) is the rate of the wave.

The constant solutions of \eqref{eq:main} are \(v=0\) and \(v=-2a_0\).
They correspond to the two possible choices of the background
\(w_\ast\). Changing that choice replaces the pair \((P,v)\) by
\begin{equation}\label{eq:shift}
 \bigl(P-2P(0),\ v+2P(0)\bigr),
\end{equation}
which again satisfies \eqref{eq:main}, as a direct substitution shows.
In the evolution equation the change of background is a Galilean shift
\(u\mapsto u+k\), \(x\mapsto x-kt\); it changes the speed relative to
the background, hence the sign of \(a_0\). The two pairs describe the
same solution \(u(x,t)\) of the same evolution equation. They coincide
exactly when \(a_0=0\). For example, the KdV soliton gives the pair
\(P=\rho^2-1\), \(v=3\operatorname{sech}^2(z/2)\). Measured from the
other background it is \(P=\rho^2+1\) and \(v-2\), which does not tend
to zero and is therefore not a rational-exponential pair.

\subsection{Dispersive, dissipative and mixed equations}

In \eqref{eq:pde} the term \(b_j\partial_x^{j+1}u\) is dispersive for
even \(j\) and dissipative for odd \(j\): for real coefficients it
changes the phase, respectively the amplitude, of a Fourier mode (with
the signs of \eqref{eq:pde}, a positive coefficient of \(u_{xx}\) is
gain). The equation is \emph{purely dispersive} if \(b_j=0\) for odd
\(j\). Then \(p\) is even, \(P\) is an even polynomial, and the equation
is invariant under \((x,t)\mapsto(-x,-t)\). The equation is \emph{purely
dissipative} if \(b_j=0\) for even \(j\). Then \(p\) is odd, \(P-P(0)\)
is an odd polynomial, and the equation is invariant under
\((x,u)\mapsto(-x,-u)\); the Nikolaevskiy equation (\(p=5\)) is an
example, see \citet{SimbawaMatthewsCox2010}. All other equations are
\emph{mixed}.

The two pure subfamilies are described by one condition. We call a monic
polynomial \(P\) of degree \(p\) \emph{reflection-symmetric} if
\(P-P(0)\) has the parity of \(\rho^p\):
\begin{equation}\label{eq:symmetric-class}
 P(-\rho)=P(\rho)\quad(p\text{ even}),\qquad
 P(\rho)+P(-\rho)=2P(0)\quad(p\text{ odd}).
\end{equation}
In the travelling-wave variable both symmetries are the reflection
\(z\mapsto-z\). If \(P\) is even and \(v(z)\) solves \eqref{eq:main}, so
does \(v(-z)\). If \(p\) is odd and \(P-P(0)\) is odd, so does
\(-v(-z)-2P(0)\); this is the reflection combined with the change of
background \eqref{eq:shift}.

\subsection{Poles, and the form of one-pole waves}\label{sec:systems}

To count waves we need to know what they look like. Three facts suffice:
a pole has order \(p\) and a fixed leading coefficient; with one pole
per period, a wave rational in \(e^z\) is a polynomial of degree \(p\)
in the logistic function \(Q\), and an elliptic wave is a combination of
\(\wp\) and its derivatives.

Throughout, a meromorphic function is single-valued and meromorphic on
the whole complex plane. The order of a pole of a solution of
\eqref{eq:main} is forced by dominant balance. A pole of order \(r\)
makes \(P(D)v\) of order \(r+p\) and \(v^2\) of order \(2r\); the two
can cancel only if \(r=p\). Comparing the leading coefficients gives, at
a pole \(z_0\),
\begin{equation}\label{eq:leadingpole}
 v(z)=\frac{c_\ast}{(z-z_0)^{p}}+O\bigl((z-z_0)^{1-p}\bigr),\qquad
 c_\ast=(-1)^{p-1}(p-1)!\,s_p,\qquad s_p=\frac{2(2p-1)!}{((p-1)!)^2}.
\end{equation}
This holds for every \(P\). The scale \(s_p\) of the normal forms of
Sections~\ref{sec:counting} and~\ref{sec:elliptic-counting} has the
first values \(s_2=12\), \(s_3=60\), \(s_4=280\), \(s_6=5544\), so that
\(c_\ast=-12,\,120,\,-1680\) for \(p=2,3,4\).

The meromorphic travelling waves in this paper are of three kinds:
rational in \(z\), rational in an exponential \(e^{\kappa z}\), and
elliptic; Theorem~\ref{thm:complete} says when there are no others. The
period group of a wave is its \emph{full} group of periods: trivial in
the first case, generated by one period in the second, a lattice in the
third. In the second case we choose \(\kappa\) so that
\(2\pi i/\kappa\) generates this group. The wave has \emph{one pole per
period} if all its poles are translates of a single pole by the periods.
A rational wave then has a single pole, and an elliptic wave one in each
period parallelogram.

Let \(v=R(t)\), \(t=e^{\kappa z}\), be a wave rational in an
exponential, with one pole per period. The rational function \(R\) has
finite limits at both \(t=0\) and \(t=\infty\). Indeed,
\(P(D)=P(\kappa t\partial_t)\) multiplies each monomial \(t^m\) by a
constant, whereas squaring doubles an extreme exponent: if \(R\) had a
pole at either end, \(R^2\) would contain an extreme Laurent monomial
that \(P(\kappa t\partial_t)R\) cannot cancel. Both limits are constant
solutions of \eqref{eq:main}. After a translation, the only pole of
\(R\) in \(0<|t|<\infty\) is at \(t=-1\), of order \(p\). For
\(\kappa=1\) such a function is a polynomial of degree \(p\) in the
logistic function \(Q=1/(1+e^{-z})=(1+\tanh(z/2))/2\), with
\(DQ=Q(1-Q)\), which has a simple pole with residue one at \(z=i\pi\)
and the limits \(0\) and \(1\). By \eqref{eq:leadingpole} the
coefficient of \(Q^p\) is \(c_\ast\); the KdV soliton above is
\(12Q(1-Q)\). Polynomial logistic representations and their use in
constructing travelling waves are developed by
\citet{Kudryashov2012,Kudryashov2013,Kudryashov2015}, and applied to
equations of the family up to \(p=6\) by
\citet{RyabovSinelshchikovKochanov2011}.

Such a wave is a \emph{pulse} if its two limits coincide and a
\emph{front} otherwise. For a real rate and a real profile these are the
solitary waves and the kinks of the evolution equation, and they are
smooth and bounded on the real axis when the pole lies off it. In
\eqref{eq:scaling} we take \(w_\ast\) to be the limit as
\(e^{\kappa z}\to0\). Then \(v\to0\) at that end, and the other limit
is \(0\) for a pulse and \(-2a_0\) for a front.

An elliptic wave whose only pole modulo its periods is at zero is a
constant plus a linear combination of \(\wp,\wp',\ldots,\wp^{(p-2)}\),
in which \(\wp^{(p-2)}\) has the coefficient \(-s_p\)
(Section~\ref{sec:elliptic-counting}). Rational waves are not counted,
but they govern what happens at infinity in both counts. A rational
solution has no nonconstant polynomial part at infinity: such a part
would have its degree doubled by \(v^2\), while the linear differential
expression cannot reach that degree. A rational wave whose only pole is
at zero is therefore
\begin{equation}\label{eq:rational-system}
 v(z)=c_\ast z^{-p}+\bigl(\text{a polynomial of degree at most }p-1
 \text{ in }z^{-1}\bigr).
\end{equation}
The simplest one is \(v=c_\ast z^{-p}\), for \(P=D^p\); it is given for
every order by \citet{Kudryashov2004}.

Conversely, a function of one of these forms that satisfies
\eqref{eq:main} is a nonconstant wave with one pole per period, because
\(c_\ast\ne0\). Substitution turns \eqref{eq:main} into finitely many
polynomial equations in the coefficients of the form and of \(P\),
without any genericity assumption. They are instances of the
principal-part construction of
\citet{DeminaKudryashov2011,DeminaKudryashovElliptic2011}; for a fixed
operator it reduces to the coefficient comparison of
\citet{DeminaKudryashov2014}, which we do not need.

\subsection{What is counted}

We count, over \(\CC\), pairs \((P,v)\) in which \(P\) ranges over
\emph{all} monic polynomials of degree \(p\) and \(v\) solves
\eqref{eq:main}. In a \emph{rational-exponential pair}, \(v\) is
rational in \(e^z\), its poles are the points of \(i\pi+2\pi i\ZZ\), and
\(v\to0\) as \(e^z\to0\). In an \emph{elliptic pair} on a lattice
\(\Lambda\), \(v\) is elliptic with period lattice \(\Lambda\), including
its scale, and its poles are the points of \(\Lambda\). Fixing the
periods removes the scale \(\kappa\) in \eqref{eq:scaling}, and placing
the pole removes translations. A wave and its mirror image are counted
as two pairs when they are distinct (Section~\ref{sec:reflection}). Both
pairs \eqref{eq:shift} of an elliptic wave are counted; for a
rational-exponential wave the condition \(v\to0\) selects one of them.
